\documentclass[11pt]{article}
\usepackage[margin=1in]{geometry}
\usepackage[T1]{fontenc}
\usepackage{lmodern}
\usepackage{amsmath,amssymb,amsthm}
\usepackage{graphicx}
\usepackage{hyperref}

\title{A Counterexample to Necessity of Almost-Homogeneous Difference Sets}
\author{Run fa\_banach\_001, agent\_lane\_04}
\date{2026-06-23}

\newtheorem*{theorem}{Theorem}
\newtheorem*{claim}{Claim}
\newcommand{\slog}{\operatorname{slog}}
\newcommand{\ellTwo}{\ell_2}
\newcommand{\R}{\mathbb R}
\newcommand{\N}{\mathbb N}

\begin{document}
\maketitle

\section*{Status}

This is a candidate counterexample packet, likely valid pending expert review.
It gives a negative answer to Question~4.3 of Margaris,
\emph{Almost bi--Lipschitz embeddings using covers of balls centred at the
origin}, arXiv:2001.02607, under the literal reading of ``linear
bi--Lipschitz embeddings into Euclidean space''.

\section*{Source Question}

Question~4.3 asks whether the condition that \(X-X\) is almost homogeneous at
\(0\) is necessary in order to have linear bi--Lipschitz embeddings into some
Euclidean space. The source question appears in the final section, immediately
after Questions~4.1 and 4.2.

\begin{center}
\includegraphics[width=\textwidth]{figures/open_problem_crop.png}
\end{center}

\section*{Idea of the Proof}

The tempting positive argument would be: a finite-dimensional image has
controlled covering numbers, so pull those covers back to \(X-X\). The flaw is
that a linear map may be bi--Lipschitz on \(X\) without being bi--Lipschitz on
the metric space \(X-X\). A pair of difference vectors need not itself be a
difference of two points of \(X\).

We exploit exactly this gap. Build \(X\) as the graph of a Lipschitz curve over
the interval. Projection onto the interval is then a linear bi--Lipschitz
embedding into \(\R\). But at scale \(h_n=2^{-n^3}\), the curve has \(2^n\)
tiny straight pieces pointing in mutually orthogonal directions. The
corresponding difference vectors lie near the origin, are mutually separated at
scale \(h_n\), and are far too numerous for any almost-homogeneous logarithmic
covering bound.

\section*{Counterexample}

\begin{theorem}
There is a compact subset \(X\) of the Hilbert space
\(\R\oplus_2\ell_2\) such that the coordinate projection
\[
  P(t,u)=t
\]
is a linear bi--Lipschitz embedding of \(X\) into \(\R\), but \(X-X\) is not
\((\alpha,\beta)\)-almost \((M,s)\)-homogeneous at \(0\) for any finite
parameters \(M\geq 1\), \(s\geq 0\), and \(\alpha,\beta\geq 0\).
\end{theorem}

\begin{proof}
Let
\[
  m_n=2^n,\qquad h_n=2^{-n^3}\qquad (n\geq 2).
\]
Since \(\sum_{n\geq 2}m_nh_n<1\), choose pairwise disjoint intervals
\[
  I_{n,k}=[a_{n,k},a_{n,k}+h_n]\subset [0,1],
  \qquad 1\leq k\leq m_n.
\]
Let \((e_{n,k})_{n\geq2,\ 1\leq k\leq m_n}\) be an orthonormal family in
\(\ell_2\). Define a measurable function \(g:[0,1]\to\ell_2\) by
\[
  g(t)=e_{n,k}\quad\text{for }t\in I_{n,k},\qquad
  g(t)=0\quad\text{otherwise},
\]
and set
\[
  f(t)=\int_0^t g(u)\,du.
\]
Then \(\|g(t)\|\leq 1\), so \(f\) is \(1\)-Lipschitz. Put
\[
  X=\{(t,f(t)):0\leq t\leq 1\}\subset \R\oplus_2\ell_2.
\]
This is compact because it is the graph of a continuous map on a compact
interval.

The coordinate projection \(P(t,u)=t\) is bounded and linear. For
\(x_t=(t,f(t))\) and \(x_u=(u,f(u))\) in \(X\),
\[
  |t-u|\leq \|x_t-x_u\|
  =\big(|t-u|^2+\|f(t)-f(u)\|^2\big)^{1/2}
  \leq \sqrt{2}\,|t-u|,
\]
because \(f\) is \(1\)-Lipschitz. Hence \(P|_X:X\to\R\) is bi--Lipschitz.

It remains to show that the difference set is not almost homogeneous at the
origin. For each \(n\geq2\) and \(1\leq k\leq m_n\), take the two endpoints of
the interval \(I_{n,k}\). Their difference in \(X\) is
\[
  z_{n,k}
  = (h_n, f(a_{n,k}+h_n)-f(a_{n,k}))
  = (h_n,h_ne_{n,k})\in X-X.
\]
Thus \(\|z_{n,k}\|=\sqrt2\,h_n\). For \(k\neq j\),
\[
  \|z_{n,k}-z_{n,j}\|
  =h_n\|e_{n,k}-e_{n,j}\|
  =\sqrt2\,h_n.
\]
Set \(r_n=2h_n\) and \(\rho_n=h_n/4\). Then all \(m_n=2^n\) points
\(\{z_{n,k}\}_{k=1}^{m_n}\) lie in \((X-X)\cap B(0,r_n)\), and no ball of
radius \(\rho_n\) can contain two of them. Consequently
\[
  N\big((X-X)\cap B(0,r_n),\rho_n\big)\geq 2^n.
\]

Suppose, toward a contradiction, that \(X-X\) were almost homogeneous at \(0\).
Then for some fixed \(M,s,\alpha,\beta\) and all \(0<\rho<r\),
\[
  N\big((X-X)\cap B(0,r),\rho\big)
  \leq M\left(\frac r\rho\right)^s\slog(r)^\alpha\slog(\rho)^\beta.
\]
Using \(r_n/\rho_n=8\), \(r_n=2^{1-n^3}\), and \(\rho_n=2^{-2-n^3}\), the
right-hand side at \((r_n,\rho_n)\) is bounded by
\[
  M8^s\,\slog(r_n)^\alpha\slog(\rho_n)^\beta
  \leq C\,n^{3(\alpha+\beta)}
\]
for a constant \(C\) depending only on the proposed parameters. This is
eventually smaller than \(2^n\), contradicting the lower bound above. Hence
\(X-X\) is not almost homogeneous at \(0\) for any finite parameters.
\end{proof}

\section*{Verification Notes}

The construction is a compact Lipschitz graph, so no completion or closure
issue is hidden. The projection is bi--Lipschitz on \(X\) with distortion at
most \(\sqrt2\). The obstruction to almost homogeneity uses a fixed covering
ratio \(r_n/\rho_n=8\), so logarithmic corrections can contribute only powers
of \(n\), while the number of separated difference vectors is \(2^n\).

The packet answers only Question~4.3 as worded. It does not settle
Question~4.1 about all doubling subsets of Banach spaces, nor Question~4.2
about finding a nonlinear almost bi--Lipschitz preprocessing map whose image
difference set is almost homogeneous at \(0\).

\section*{Novelty Check}

The cheap run indexes were searched for arXiv:2001.02607, the title, Margaris,
Olson--Robinson, \texttt{almost homogeneous at 0}, and
\texttt{linear bi-Lipschitz}. Local source search and a targeted web search for
the exact Question~4.3 phrase and close variants found the source paper but no
later answer. This is a bounded novelty check, not a comprehensive literature
survey.

\section*{Human Review Recommendation}

Review the interpretation of ``linear bi--Lipschitz embedding'' in
Question~4.3. If it means a bounded linear operator on the ambient Banach space
whose restriction to \(X\) is bi--Lipschitz, the counterexample gives a
negative answer.

\begin{thebibliography}{9}
\bibitem{Margaris}
A.~Margaris,
\emph{Almost bi--Lipschitz embeddings using covers of balls centred at the
origin},
arXiv:2001.02607.
\end{thebibliography}

\end{document}
